%%==========================================================================%%
%% background.tex -- "Background & Summary".                                  %%
%% Motivation + contents + reuse value, citing the prior publication that     %%
%% used the data. No subjective novelty, impact, or utility claims.           %%
%% American spelling. Cohort numbers are [TBD] or approximate.                %%
%%==========================================================================%%

\section*{Background \& Summary}\label{sec:background}

Whole-body magnetic resonance imaging (MRI) provides soft-tissue-rich anatomical coverage of the body without the ionizing radiation of computed tomography (CT).
In pediatric oncology, and in lymphoma in particular, this property is well suited to diagnosis, staging, and treatment monitoring, where limiting cumulative radiation across repeated examinations is a primary concern~\cite{GalloBernal_2025,Kurch_2021}.
Children are more sensitive to ionizing radiation than adults, which is one reason MRI is favored for the anatomical component of pediatric whole-body imaging~\cite{Brenner_2007}.
Whole-body MRI is increasingly used alongside or in place of CT-based protocols for these pediatric applications~\cite{Masselli_2025}.

Image-based artificial intelligence increasingly supports these clinical tasks, and a common prerequisite is the reliable identification of anatomy across the whole body, against which functional or pathological findings can be localized and scored.
Progress in automated segmentation has been driven by large public imaging resources, but these remain overwhelmingly adult and dominated by CT~\cite{Wasserthal_2023}.
The adult MRI counterpart, an open sequence-independent multi-structure segmentation tool, has recently become available for adult anatomy~\cite{Wasserthal_2025}.
Adult-trained models do not transfer cleanly to children.
Pediatric anatomy changes substantially across development, so models calibrated on adults behave unpredictably on pediatric inputs~\cite{Chatterjee_2025}, and adult-trained segmentation models degrade across pediatric tasks, with the largest losses in the youngest patients~\cite{Laborie_2026}.

Pediatric segmentation resources are beginning to emerge, but each leaves part of the whole-body MRI problem uncovered.
Pediatric abdominal organ segmentation has been demonstrated on CT through transfer learning~\cite{Somasundaram_2024}, and a recent pediatric effort segments thoracic and abdominal viscera on CT and MRI~\cite{Eicke_2026}, but neither covers the skeleton in a whole-body MRI setting.
Adult full-body and torso MRI multi-structure segmentation, including skeletal structures, has been released for large adult biobanks~\cite{Graf_2026}, and whole-body semantic anatomy segmentation has been demonstrated on adult PET/CT~\cite{Sundar_2022}, but these are adult resources.
Public pediatric whole-body MRI datasets with manual segmentations covering the major skeletal structures together with the principal organs remain scarce.

We present a pediatric whole-body MRI dataset assembled from a single-center lymphoma cohort, with a subset of examinations carrying manual multi-structure normal-anatomy segmentations.
The collection comprises approximately [TBD] whole-body MRI examinations.
A subset of approximately [TBD] examinations carries segmentation labels covering major skeletal structures and the principal abdominal and thoracic organs, contoured in 3D Slicer.
An overview of the cohort and the imaging is shown in Figure~\ref{fig:cohort}.
The same source patients and co-acquired imaging also underlie a companion descriptor that releases the positron emission tomography (PET) examinations with tumor-lesion segmentations.
The two descriptors are split by labeled task, normal-anatomy segmentation here and lesion segmentation in the companion, rather than by raw modality, and they share a portion of the underlying cohort [TBD: cross-reference the companion descriptor and its data citation once its Digital Object Identifier is minted].

The segmentation subset of this dataset was used to develop and report a pediatric whole-body MRI anatomy-segmentation model in a separate study~\cite{[TBD: companion findings publication, cite once available]}.
We anticipate that the dataset will support the development of pediatric whole-body anatomical segmentation methods, adult-to-pediatric transfer learning, and the use of normal-anatomy reference labels for downstream lesion-detection and staging research on the same imaging.
By providing pediatric whole-body MRI with consistent organization and labeling, it complements existing adult and CT-dominated resources.

%%--------------------------------------------------------------------------%%
%% Figure 1 -- cohort and example overview. Placeholder, real image pending. %%
%% Distribution-panel numbers are [TBD].                                     %%
%%--------------------------------------------------------------------------%%
\begin{figure}[t]
\centering
\fbox{\parbox[c][6cm][c]{0.85\textwidth}{\centering\itshape
[Figure 1 placeholder: (a) representative pediatric coronal whole-body MRI
with the multi-structure anatomy-segmentation overlay (skeleton and organs).
(b) Composition of the dataset across the full MRI cohort and the
anatomy-segmentation subset. (c) Age and sex distribution of the
anatomy-segmentation subset. Real image pending.]}}
\caption{\textbf{Overview of the pediatric whole-body MRI cohort and the anatomy-segmentation subset.}
(a) A representative whole-body MRI examination shown in the coronal plane with the multi-structure normal-anatomy segmentation overlay covering skeletal structures and major organs.
(b) Composition of the dataset across the full MRI cohort and the anatomy-segmentation subset (counts [TBD]).
(c) Distribution of patient age and sex for the anatomy-segmentation subset (counts [TBD]).}
\label{fig:cohort}
\end{figure}
